\documentclass[11pt,a4paper]{article}

\usepackage[a4paper,margin=2.2cm]{geometry}
\usepackage[utf8]{inputenc}
\usepackage[T1]{fontenc}
\usepackage{booktabs}
\usepackage{xcolor}
\usepackage{hyperref}
\usepackage{amsmath}

\hypersetup{
    colorlinks=true,
    linkcolor=black,
    urlcolor=blue
}

\setlength{\parindent}{0pt}
\setlength{\parskip}{0.8em}
\setlength{\emergencystretch}{2em}

\begin{document}

\begin{center}
    {\Large Informationsvisualisierung -- Übung 9: Explizite Baumdarstellungen}
\end{center}

\section*{Lauffähiger Code}

\href{https://github.com/yaniktfl/iv-repo}{https://github.com/yaniktfl/iv-repo}

\section*{Ausführung}

\begin{verbatim}
cd 9
elm reactor
\end{verbatim}

Anschließend im Browser \texttt{http://localhost:8000} öffnen und dort
\texttt{Main.elm} auswählen.

\section*{Datenstruktur und Layout}

Jeder Knoten wird einheitlich als \texttt{Node} modelliert; eine eindeutige
ID wird nach dem Aufbau über \texttt{Tree.indexedMap} vergeben (für das
Hervorheben beim Hovern):

\begin{verbatim}
type alias Node =
    { id : Int, label : String, value : Maybe Int }

assignIds : Tree Node -> Tree Node
assignIds =
    Tree.indexedMap (\id node -> { node | id = id })
\end{verbatim}

Das Layout berechnet \texttt{Hierarchy.tidy}. Es liefert für jeden Knoten die
Koordinaten \texttt{x}/\texttt{y} sowie \texttt{width}/\texttt{height} zurück.
Die eigentliche Zeichenaufgabe -- Knoten als \texttt{circle}, Kanten als
\texttt{line} -- wird selbst implementiert:

\begin{verbatim}
positionedTree = Hierarchy.tidy layoutAttrs sourceTree
positionedNodes = Tree.toList positionedTree
links           = Tree.links positionedTree

linkElement ( parent, child ) =
    line [ x1 parent.x, y1 parent.y, x2 child.x, y2 child.y ] []
\end{verbatim}

\section*{Aufgabe 9.1: Binärer und geordneter Baum}

Beide Bäume werden mit \texttt{Tree.tree} bzw. dem Helfer \texttt{leaf}
(für Blätter) direkt im Code aufgebaut. Der binäre Baum bildet Folie~405 nach,
der geordnete Baum Folie~412:

\begin{verbatim}
orderedTree =
    assignIds
        (Tree.tree (mkNode "root" Nothing)
            [ Tree.tree (mkNode "group-a" Nothing)
                [ leaf "a1", leaf "a2", ... ]
            , leaf "single-b"
            , ...
            ] )
\end{verbatim}

Für diese kleinen Bäume wird das Layout über
\texttt{Hierarchy.size 1320 520} fest auf die Zeichenfläche skaliert. Knoten
erhalten den Radius~20, Labels werden hier nicht angezeigt (nur beim Hovern).

\section*{Aufgabe 9.2: flare.json laden und darstellen}

Die Datei wird beim Start per \texttt{Http.get} geladen und über
\texttt{Cmd.batch} parallel zur Länderdatei angefordert. Ein rekursiver
Decoder überführt das geschachtelte JSON in einen \texttt{Tree Node}. Ein
Knoten hat entweder \texttt{children} (innerer Knoten) oder einen
\texttt{value} (Blatt):

\begin{verbatim}
flareDecoder : Decoder (Tree Node)
flareDecoder =
    Decode.map3
        (\name value children ->
            Tree.tree (mkNode name value) (Maybe.withDefault [] children))
        (Decode.field "name" Decode.string)
        (Decode.maybe (Decode.field "value" Decode.int))
        (Decode.maybe
            (Decode.field "children"
                (Decode.list (Decode.lazy (\_ -> flareDecoder)))))
\end{verbatim}

\texttt{Decode.lazy} ist nötig, da der Decoder sich für die Kinder selbst
aufruft. \texttt{Decode.maybe} fängt fehlende Felder ab, sodass derselbe
Decoder für innere Knoten und Blätter funktioniert.

\subsection*{Labels um 90 Grad drehen}

Da der Baum sehr breit ist, würden waagerechte Labels überlappen. Sie werden
deshalb wie gefordert um 90$^\circ$ um ihren eigenen Ankerpunkt gedreht,
sodass sie senkrecht und gut lesbar stehen:

\begin{verbatim}
rotation =
    if config.rotateNodeLabels then
        [ transform [ Rotate 90 labelX labelY ] ]
    else
        []
\end{verbatim}

\section*{Aufgabe 9.3: Eigener Decoder für die Länderhierarchie}

\texttt{countryHierarchy.json} hat ein anderes Format: Der Name steht
verschachtelt unter \texttt{data.id}, und es gibt keine Werte an den Blättern.
Dafür wird ein eigener Decoder geschrieben, der den Namen über
\texttt{Decode.at} aus dem geschachtelten Feld liest:

\begin{verbatim}
countryDecoder : Decoder (Tree Node)
countryDecoder =
    Decode.map2
        (\name children ->
            Tree.tree (mkNode name Nothing) (Maybe.withDefault [] children))
        (Decode.at [ "data", "id" ] Decode.string)
        (Decode.maybe
            (Decode.field "children"
                (Decode.list (Decode.lazy (\_ -> countryDecoder)))))
\end{verbatim}

Der resultierende Baum (World $\to$ Kontinent $\to$ Subregion $\to$ Land) wird
mit demselben \texttt{treeView} dargestellt wie flare.

\section*{Lesbarkeit großer Bäume}

flare ($\approx$250 Knoten) und die Länderhierarchie (261 Knoten) sind deutlich
größer als die Bäume aus 9.1. Zunächst wurden alle vier Bäume gleich behandelt
und mit \texttt{Hierarchy.size 1700 850} gezeichnet. Bei den großen Bäumen
funktioniert das nicht, weil die rund 220 Blätter von flare auf eine feste
Breite von 1700\,px gequetscht werden und Knoten wie Label dadurch übereinander
liegen. Zusätzlich verkleinerte das SVG über \texttt{width:\,100\%} alles noch
auf die Seitenbreite.

Der Grund liegt darin, dass \texttt{Hierarchy.size} und
\texttt{Hierarchy.nodeSize} unterschiedlich arbeiten. \texttt{size} skaliert
das fertige Layout in eine feste Box, was bei kleinen Bäumen passt, deren
Größe man kennt. \texttt{nodeSize} dagegen gibt jedem Knoten einen festen Platz
und lässt das Layout entsprechend wachsen. Für hunderte Knoten ist nur das
Zweite sinnvoll. Die Knoten oder die Schrift einfach kleiner zu machen, hätte
das Überlappen nur verschoben, denn die Ursache ist das Einpassen in eine feste
Fläche.

Deshalb werden die beiden Fälle über das Feld
\texttt{fixedSize : Maybe (Float, Float)} getrennt. Kleine Bäume behalten ihre
feste Fläche, große Bäume lassen \texttt{Hierarchy.size} weg und wachsen über
\texttt{nodeSize} und die Abstände natürlich.

\begin{verbatim}
layoutAttrs =
    (case config.fixedSize of
        Just ( w, h ) -> [ Hierarchy.size w h ]
        Nothing       -> [])
    ++ [ Hierarchy.nodeSize (\_ -> ( r * 2, r * 2 ))
       , Hierarchy.parentChildMargin config.parentChildMargin
       , Hierarchy.peerMargin config.peerMargin
       , Hierarchy.layered ]
\end{verbatim}

Die \texttt{viewBox} wird anschließend aus den tatsächlichen Knoten-Grenzen
(Minimum/Maximum von \texttt{x}/\texttt{y}) berechnet, mit zusätzlichem Platz
unten für die gedrehten Labels. Das SVG wird in seiner echten Pixelgröße in
einem scrollbaren Container gerendert, statt auf Seitenbreite herunterskaliert
zu werden -- so bleiben die Label lesbar und man scrollt durch den breiten Baum.

\begin{center}
\begin{tabular}{lll}
\toprule
Baum & \texttt{Hierarchy.size} & Darstellung \\
\midrule
9.1 binär / geordnet & $1320 \times 520$ (fest) & skaliert in feste Fläche \\
9.2 flare / 9.3 Länder & -- (natürlich) & echte Größe, scrollbar, Labels 90$^\circ$ \\
\bottomrule
\end{tabular}
\end{center}

\end{document}
